\section{The intrinsic idempotent core}\label{sec:intrinsic}
\subsection{Idempotency identifies an intrinsic core}\label{gen:sec:core}

A positive-arity term operation is \emph{idempotent} when
$t(S,\ldots,S)=S$ for every nonempty $S$. This semantic condition
identifies a smaller syntax. For a variable occurrence, compose all
role permutations on its path from the root, in the order in which
they act on its value.

\begin{theorem}[Exact idempotent elimination]\label{gen:thm:core}
A term is idempotent if and only if it has no winner leaves and every
variable occurrence has identity accumulated permutation. Consequently
the idempotent part of the full term clone is exactly the clone generated
by the $K$ binary owned choices, without constants or internal switches.
Recognition and extraction of such a term take polynomial time in the
explicit input description.
\end{theorem}
\begin{proof}
On the all-ambiguous cell, track a constant set $C$ and the
variable-occurrence permutations. If any winner leaf occurs, $C$ is
nonempty: it starts nonempty at that leaf, is transported by switches,
survives unions, and a collapse at an owned node replaces it by a
nonempty owner singleton. A node with no constant contribution cannot
collapse on this cell, since every live child is ambiguous. Thus a term
containing a winner leaf has some fixed nonempty $C$ in every ambiguous
diagonal output. But $\bigcap_{|S|\geq2}S=\varnothing$ for $K\geq3$,
so the term cannot be idempotent.

Without winner leaves, every intermediate value on this cell is
ambiguous. Every choice is union. The diagonal value is therefore the
union of $\pi S$ over all variable occurrences. Idempotency forces each
$\pi S\subseteq S$. If $\pi(i)\ne i$, a pair containing $i$ and omitting
$\pi(i)$ contradicts that inclusion. Every accumulated permutation is
identity.

Transport switches through choices using
$\pi(u\star_r v)=\pi u\star_{\pi(r)}\pi v$ and compose them.
All leaf permutations disappear. The remaining constant-free owned term
is idempotent because its primitive choices are. A traversal accumulating
permutations both checks the criterion and transports the owners,
without expanding the tree.
\end{proof}

The no-switching calculus thus contains \emph{every} idempotent definable
operation. Internal switching still enlarges the complete clone and
changes its congruences, as \cref{view:thm:sym} shows.

\begin{corollary}[Semilattice trace]\label{gen:cor:trace}
On all-ambiguous inputs, every idempotent term is union over a fixed
nonempty subset of its arguments. These restrictions form exactly the
semilattice term clone, and restriction is a surjective clone
homomorphism. There is no near-unanimity term of any finite arity, and
no Mal'cev term. Every variable occurring in a core term is semantically
essential, so every valid core identity has the same variables on both
sides.
\end{corollary}
\begin{proof}
In the constant-free core all intermediate ambiguous values are
ambiguous; the result is union of the occurring variables. Distinct
supports are separated by a pair and $I$. Composition unions supports,
and every nonempty support is realized by derived union. A nonempty
join cannot satisfy near-unanimity: put the larger of those two values
at a selected argument and the smaller at all others. The Mal'cev
identities would require support both in the first argument alone and
in the third alone. Operations satisfying either collection of
identities are already idempotent. The same two-value test detects the
presence of each argument, proving essentiality. Regularity of identities
is the standard two-element semilattice observation: interpret every
primitive choice as join on $\{P,I\}$, for a proper ambiguous pair $P$.
An identity must then have the same variable support on both sides.
This is the usual regular-identity setting of semilattice sums
\cite{Plonka1967}.
\end{proof}

\begin{corollary}[Intrinsic player views]\label{gen:cor:views}
Every idempotent operation in the full role-switching language acts
compositionally on each of the $K+1$-state player views
\[
 q_k(S)=
 \begin{cases}
 \pv,&k\in S,\\
 a_i,&S=\{i\},\ i\ne k,\\
 \nv,&\text{otherwise}.
 \end{cases}
\]
These give a subdirect representation by $K$ factors. The family is
irredundant: none of its player views can be omitted.
\end{corollary}
\begin{proof}
Each $q_k$ is a congruence quotient for every owned choice by \cref{view:thm:view}. Theorem~\ref{gen:thm:core} extends this to all
idempotent operations, even when their original definitions contain
internal switches. All membership bits are recovered from the $\pv$
states, so the family separates values. Each view is onto. Omitting
$q_k$ identifies $I\setminus\{k\}$ and $I$, which every other player
views as $\pv$.
\end{proof}

The decomposition that fails for unrestricted internal role switching
therefore holds intrinsically for all idempotent definable operations.
The representation is subdirect, not the full Cartesian product.

\subsubsection{An optimal logarithmic test set}

A fixed domain $D\subseteq\Pow(I)$ \emph{certifies idempotency} when,
for terms of every arity, agreement $t(S,\ldots,S)=S$ for $S\in D$
implies idempotency. This tests a specific diagonal law, rather than
all term equations.

\begin{theorem}[Exact ambiguous certificates]\label{gen:thm:idem-domain}
For a domain consisting of ambiguous values, certification holds exactly
when
\[
 \bigcap D=\varnothing,\qquad
 (\mathbf1_{i\in S})_{S\in D}\text{ are distinct for all }i\in I.
\]
The minimum size of a certifying domain, even allowing singleton test
values, is three for $K=3$, and $\lceil\log_2(K+1)\rceil$ for $K\geq4$.
Both minima are attained using ambiguous values alone.
\end{theorem}
\begin{proof}
If $k\in\bigcap D$, the nonidentity term $x\cup\{k\}$ fixes every test.
If two roles have the same incidence vector, their transposition fixes
every test. These necessary conditions apply also to singleton tests.

Conversely suppose every test is ambiguous and the two conditions hold.
An agreeing term's constant contribution lies in $\bigcap D$, so is
empty. Hence it has no winner leaves. Each occurrence permutation
satisfies $\pi S\subseteq S$ on each test, and finiteness makes these
equalities. A permutation fixing each test can only exchange roles with
identical incidence vectors. It is identity. Apply Theorem~\ref{gen:thm:core}.

With $m$ tests the conditions require $K$ distinct binary words of
length $m$, excluding the all-one word. Thus $K\leq2^m-1$. For $K\geq4$
set $m=\lceil\log_2(K+1)\rceil$. Start the role signatures with the $m$
distinct two-bit words supported on a cycle on $m$ bits. Each bit occurs
twice, making every test ambiguous. Add distinct words other than the
all-one word until there are $K$ signatures.

For $K=3$, two tests satisfying the necessary conditions have signatures
$00,01,10$, hence are two distinct singletons. The reduct's ambiguity
adder from \cref{game:thm:unary} fixes every singleton but adds $c$ to
ambiguous inputs, so those tests cannot certify idempotency.
Three pairs certify it.
\end{proof}

At $K=7$, three ambiguous diagonal tests suffice, compared with $28$
values for testing every full-signature equation. The improvement is
specific to idempotency.

\subsection{Separation and deterministic atomic expressivity}\label{gen:sec:sep}

Fix input environments $v_a$ and pairwise disjoint target regions
$W_a\subseteq I$. A source occurrence is a winner leaf, or a variable
leaf together with its accumulated permutation.

\begin{theorem}[One common source]\label{gen:thm:source}
If $t(v_a)\subseteq W_a$ for every $a$, one source occurrence of $t$
already has its value inside $W_a$ in every environment simultaneously.
If there are at least two target regions, it is a variable occurrence.
\end{theorem}
\begin{proof}
At an owned node with owner $r$, at most one target contains $r$.
In every other environment the output cannot collapse to $\{r\}$;
it is union, so both children lie in the target. In the possible
exceptional environment, either both children lie in the target or a
child is the absorbing owner singleton and lies in it. Choose that child:
it works in all environments. Descend to a leaf. At a switch pull the
targets back along the permutation, preserving disjointness. A winner
leaf has the same nonempty value in all environments and cannot fit
two disjoint targets.
\end{proof}

Distinct singleton outputs at any number of tests are thus jointly
explained by one occurrence. It cannot change between tests.
\texttt{Separation.lean} proves this for arbitrary index types and target
predicates, with internal permutations explicit in the syntax.

Call a term \emph{atom-preserving} if it returns a singleton on every
singleton-input valuation. It induces a function $f:I^n\to I$.

\begin{theorem}[Exact deterministic atomic clone]\label{gen:thm:S2}
At every positive arity, the functions induced by atom-preserving terms
are exactly the essentially unary or at most two-valued functions.
This is the known clone $S_2$.
\end{theorem}
\begin{proof}
Let $\iota_3\subseteq I^3$ be the relation of triples with a repeated
entry. Three input environments having this property in every
coordinate cannot have three distinct singleton outputs:
Theorem~\ref{gen:thm:source} would give one permutation sending an input
triple with repetition to three distinct roles. Every induced function
therefore preserves $\iota_3$.

Use the established presentation $\operatorname{Pol}(\iota_3)=S_2$,
recorded by Szendrei\cite[Proposition~2.2]{Szendrei2012} in the
S\l upecki--Burle classification. It applies for every finite $K\geq3$.
This is a use of that classification, not a new classification of $S_2$.

Conversely every unary map $h:I\to I$ has an atomic realization.
The term $x\star_a(a\ b)(x)$ merges atomic $b$ into $a$, fixing the
others. Merge each $h$-fiber into a representative, then permute the
representatives to the prescribed images.

For distinct output roles $a,b$, tests of input labels are unary maps
into $\{a,b\}$. On these atomic values, $\star_a$ is Boolean OR and
$\star_b$ Boolean AND, with truth $a$. Finite DNF of label tests
realizes every $I^n\to\{a,b\}$. Constants cover the one-valued case.
\end{proof}

\begin{corollary}[Idempotent atomic determinacy]
An idempotent, atom-preserving term projects one input on all atomic
valuations. A genuinely multiargument rule using at least three output
roles must return ambiguity somewhere on atomic inputs.
\end{corollary}
\begin{proof}
Idempotency makes its atomic function surjective. As $K\geq3$,
Theorem~\ref{gen:thm:S2} makes it essentially unary, and the diagonal law
makes that unary map identity.
\end{proof}

The exact number of induced $n$-ary atomic functions is
\[
 \binom K2(2^{K^n}-2)+K+
 n\left(K^K-\binom K2(2^K-2)-K\right).
\]
The first part counts functions with at most two values. The remaining
essentially unary functions have a unique essential argument.
For $K=3,n=2$ the count is $1545$. This classifies deterministic atomic
behavior, leaving mixed and ambiguous valuations available for richer
operations.
